Large pre-trained models(LPTMS) such as BERT \citep{devlin2019bert}, GPT \citep{brown2020language}, CLIP \citep{radford2021learning}, DINO \citep{caron2021emerging}, and SigLIP \citep{zhai2023sigmoid} have become foundational across vision, language, and multimodal tasks. Their learned representations transfer effectively to a wide range of downstream applications, from image classification and retrieval to visual question answering and zero-shot recognition \citep{bommasani2021opportunities}. \textcolor{red}{When these models are deployed in specialized domains such as satellite imagery, medical imaging, or fine-grained recognition, their general-purpose representations often fall short, necessitating domain adaptation.} A rich line of work addresses this through parameter-efficient fine-tuning: LoRA \citep{hu2022lora} introduces low-rank updates to pretrained weights, while prompt-based methods such as CoOp \citep{zhou2022learning}, CoCoOp \citep{zhou2022conditional}, and MaPLe \citep{khattak2023maple} learn task-specific context tokens in the text or both modalities of vision-language models. Beyond few-shot prompt tuning, adapter-based approaches \citep{gao2024clipadapter} and visual prompt tuning \citep{jia2022visual} offer complementary strategies that modify intermediate representations or input-level tokens while keeping the backbone frozen. In the unsupervised setting, recent methods leverage CLIP's zero-shot priors for source-free domain adaptation \citep{singha2023ad}, and test-time prompt tuning \citep{shu2022tpt} adapts prompts on the fly using entropy minimization over unlabeled target data. \textcolor{red}{Despite the effectiveness of these adaptation strategies, how they reshape the model's internal representations, and whether prior interpretability analyses remain valid after adaptation, is largely unexplored.}

As these systems are increasingly deployed in high-stakes settings such as medical diagnosis \citep{singhal2023large}, autonomous driving \citep{cui2024survey}, and legal reasoning \citep{guha2024legalbench}, \textcolor{red}{understanding what they represent internally is no longer optional} \citep{rauker2023transparent,kim2018interpretability}. Mechanistic interpretability offers one such route. Sparse autoencoders (SAEs), which decompose polysemantic neural activations into sparse, monosemantic features, have emerged as a principled tool for resolving superposition in neural networks \citep{cunningham2024sparse, rajamanoharan2024improving}. Following their success in language models, recent work has extended SAEs to vision transformers \citep{gandelsman2024interpreting, gao2025scaling}.

\textcolor{red}{A neuron-level analysis of adapted models can reveal whether the features learned during pretraining are preserved, repurposed, or replaced during adaptation, directly informing how much trust we can place in a fine-tuned model's behavior on the target domain} \citep{durrani2020analyzing, csiszarik2021similarity}. Understanding this internal reorganization is essential: \textcolor{red}{if adaptation merely remaps existing features to new categories, the model's reliability on the target domain rests on the quality of that remapping; if adaptation induces new feature representations, we need tools to verify their faithfulness to the target distribution} rather than inheriting assumptions from the base model's explanations.

A growing body of work trains SAEs on a base pretrained model and reuses them as fixed analysis tools: Rao et al. \citep{rao2024discover} use CLIP SAE dictionaries as task-agnostic concept bottlenecks across downstream tasks, Karvonen et al. \citep{karvonen2025saevision} apply a single SAE trained on frozen ViT activations for causal edits across classification and segmentation without retraining, and Pach et al. \citep{pach2025monosemantic} show that SAE interventions on CLIP's vision encoder can steer multimodal LLM outputs without modifying the language model. Most directly relevant, Lim et al. \citep{lim2024patchsae} apply a base-model SAE to MaPLe-adapted CLIP and \textcolor{red}{find that prompt-based adaptation largely remaps existing concepts rather than learning new ones. However, their analysis is restricted to prompt tuning on near-domain benchmarks where the base SAE's dictionary remains a reasonable fit.} \textcolor{red}{We show that this assumption breaks down under stronger adaptation methods (LoRA, DoRA, TAP, ERM, DRO, L2C) and larger domain shifts (medical imaging, satellite imagery, textures), where base SAE features lose both reconstruction fidelity and causal faithfulness to the adapted model's behavior, necessitating domain-adapted SAEs and new metrics to quantify the degradation.}

\textcolor{red}{To ensure our findings generalize across model architectures, adaptation strategies, and domain characteristics, we design our experimental setup along three axes.} First, we study four vision encoders that span distinct pretraining objectives: CLIP \citep{radford2021learning} and SigLIP~2 \citep{tschannen2025siglip2}, which learn visual representations through language supervision via contrastive and sigmoid losses respectively; ALIGN \citep{jia2021scaling}, which scales vision-language pretraining with noisy web data; and DINOv2 \citep{oquab2024dinov2}, a purely self-supervised vision model with no language signal. \textcolor{red}{This selection allows us to disentangle the role of language grounding from visual feature quality when studying how adaptation affects interpretability.}

Second, we consider three families of adaptation strategies that span a spectrum of how aggressively the model is modified. Parameter-efficient fine-tuning via LoRA \citep{hu2022lora} introduces low-rank weight perturbations directly into transformer layers. Prompt-based methods, including MaPLe \citep{khattak2023maple} and TAP \citep{ding2025tap}, modify model behavior by learning input-level context tokens while leaving all backbone weights frozen. For datasets exhibiting distributional shift without labeled target data, we additionally evaluate models trained with empirical risk minimization (ERM), distributionally robust optimization (DRO) \citep{sagawa2020distributionally}, and the zero-shot transfer baseline of CLIP (ClipOOD), \textcolor{red}{enabling us to study whether SAE interpretability degrades differently under supervised fine-tuning versus robustness-oriented training.}

Third, our dataset selection is designed to systematically vary domain distance from the ImageNet-centric pretraining distribution. We organize 25 datasets into three tiers. The first tier comprises domain generalization benchmarks with explicit distribution shifts: PACS, Mini-DomainNet, Office-Home, and VLCS, which test cross-domain transfer across art styles, sketch, and photo domains; Cityscapes-KITTI and iWildCam, which introduce environmental and geographic shift in scene understanding and wildlife recognition; and Camelyon17 \citep{bandi2018detection}, which presents histopathology slides with hospital-level distributional variation. The second tier includes 16 established few-shot and transfer learning benchmarks spanning fine-grained recognition (StanfordCars, FGVCAircraft, Flowers102, OxfordPets), scene and action understanding (SUN397, UCF101), specialized domains far from natural images (EuroSAT for satellite imagery, MedMNIST and CheXpert for medical imaging, DTD for textures), and ImageNet variants (ImageNet, ImageNetV2, ImageNet-Sketch, ImageNet-A) that provide controlled covariate shift within the pretraining domain. \textcolor{red}{This tiered design enables us to test whether SAE degradation is a smooth function of domain distance or exhibits sharp phase transitions at critical levels of distributional shift.}

To study this, we use the SAE lens and explicitly position our analysis among multiple SAE variants: ReLU SAEs \citep{cunningham2024sparse}, Top-k SAEs \citep{gao2025scaling}, Gated SAEs \citep{rajamanoharan2024improving}, and Matryoshka SAEs \citep{zaremba2025msae, bussmann2025matryoshka}. \textcolor{red}{We focus on PatchSAE \citep{lim2024patchsae} because it operates over patch tokens and therefore supports both global and local concept analysis in images.} This allows us to compare adaptation methods not only by end-task accuracy, but \textcolor{red}{by how they reshape concept-level structure across datasets.}


% \textcolor{blue}{
% \begin{itemize}
%     \item \textbf{DA axis (Domain Adaptation).} Covered by ERM, DRO, and ClipOOD on PACS, Mini-DomainNet, Office-Home, VLCS, Cityscapes-KITTI, iWildCam, and Camelyon17. Basis papers: DomainBed~\citep{gulrajani2021domainbed}, WILDS~\citep{koh2021wilds}, DRO~\citep{sagawa2020distributionally}.
%     \item \textbf{DG axis (Domain Generalization / OOD robustness).} Covered by ImageNet-A, ImageNet-R, ImageNet-Sketch, and ImageNetV2 evaluated with adapted models. Basis: the OOD robustness evaluation protocol from CoOp~\citep{zhou2022learning}, CoCoOp~\citep{zhou2022conditional}, MaPLe~\citep{khattak2023maple}, and TAP~\citep{ding2025tap}, which all report on these four OOD benchmarks.
%     \item \textbf{LoRA / Prompt-tuning axis (Few-shot adaptation).} Covered by LoRA, MaPLe, and TAP on the 11-dataset PatchSAE/CoOp benchmark (Caltech101 through SUN397), plus MedMNIST, CheXpert, and EuroSAT. Basis papers: PatchSAE~\citep{lim2024patchsae}, TAP~\citep{ding2025tap}, MaPLe~\citep{khattak2023maple}, PromptSRC~\citep{khattak2023promptsrc}, ProLIP~\citep{prolip2026}.
%     \item \textbf{OOD / Distribution shift axis (Interpretability under shift).} Covered by cross-domain shift datasets (EuroSAT, MedMNIST, CheXpert) and within-domain shift datasets (ImageNet variants, DTD, UCF101).
% \end{itemize}
% }